% document class
\documentclass{book}

% list of packages
\usepackage[left=2cm,right=2cm,top=2cm,bottom=2cm]{geometry} % document size

\usepackage[utf8]{inputenc}
\usepackage[T2A]{fontenc} %russian language encoding

\usepackage{textcomp}
\usepackage{lmodern}
\usepackage{amsmath}
\usepackage{hyperref}
\usepackage{xcolor}
\usepackage{graphicx}
\usepackage{listings}
\usepackage{listingsutf8}
\usepackage{dcolumn}   % needed for some tables
\usepackage{bm}        
\usepackage{amssymb} 
\usepackage{fancyvrb}
\usepackage{fvextra}
%\usepackage{minted}

%\setcounter{secnumdepth}{3}
%\setcounter{tocdepth}{3}
%\usepackage{eso-pic}
\usepackage{array}
%\usepackage{rotating}
\usepackage{multicol}
\usepackage{setspace}

%russian language support
\usepackage[english,russian]{babel}
%starts at the last language on the list
%select language switches the language


\lstset{breaklines=true}

% The following is a list of styles for use with listings

\lstdefinestyle{codepy}{
    backgroundcolor=\color{white},   
    commentstyle=\color{green!50!black},
    keywordstyle=\color{blue},
    numberstyle=\tiny\color{gray},
    stringstyle=\color{purple},
    basicstyle=\ttfamily\footnotesize,
    breakatwhitespace=false,         
    breaklines=true,                 
    captionpos=b,                    
    keepspaces=true,                 
    numbers=left,                    
    numbersep=5pt,                  
    showspaces=false,                
    showstringspaces=false,
    showtabs=false,                  
    tabsize=2
}

\lstdefinestyle{codejupytr}{
    backgroundcolor=\color{gray!10},   
    commentstyle=\color{green!50!black},
    keywordstyle=\color{blue},
    numberstyle=\tiny\color{gray},
    stringstyle=\color{purple},
    basicstyle=\ttfamily\footnotesize,
    breakatwhitespace=false,         
    breaklines=true,                 
    captionpos=b,                    
    keepspaces=true,                 
    numbers=left,                    
    numbersep=5pt,                  
    showspaces=false,                
    showstringspaces=false,
    showtabs=false,                  
    tabsize=2
}

\lstdefinestyle{output}{
    backgroundcolor=\color{gray!10},   
    commentstyle=\color{black},
    keywordstyle=\color{black},
    numberstyle=\tiny\color{gray},
    stringstyle=\color{black},
    basicstyle=\ttfamily\footnotesize,
    breakatwhitespace=false,         
    breaklines=true,                 
    captionpos=b,                    
    keepspaces=true,                 
    numbers=left,                    
    numbersep=5pt,                  
    showspaces=false,                
    showstringspaces=false,
    showtabs=false,                  
    tabsize=2
}


\makeatletter


\graphicspath{{figs/}} % Directory in which figures are stored


%
% list of custom commands
\newcommand{\be}{\begin{equation}}
\newcommand{\ee}{\end{equation}}
\newcommand{\bea}{\begin{eqnarray}}
\newcommand{\eea}{\end{eqnarray}}
\renewcommand{\(}{\left(}
\renewcommand{\)}{\right)}

\newcommand{\rut}{\selectlanguage{russian}}
\newcommand{\ent}{\selectlanguage{english}}


%\newcommand{\be}{\begin{equation}}
%\rnc{\d}{\mathrm{d}}
%\nc{\D}{\partial}
%\nc{\nn}{\nonumber}
%\rnc{\inf}{\infty}
%\nc{\vphi}{\varphi}
%\nc{\half}{\frac{1}{2}}


% color commands are used to highlight text temporarily, for editing
\newcommand{\nadd}[1]{\textcolor{magenta}{#1}}
\newcommand{\ncomm}[1]{\textcolor{magenta}{\bf{[#1]}}}
\renewcommand{\Re}{\mathrm{Re}}
\renewcommand{\Im}{\mathrm{Im}}




\makeatother

\begin{document}

% title page
\title{Исследование влияния электрон-электронного взаимодействия на
незатухающий ток в квантовых проводах}

\author{Гоулд Элизабет Сара\\
Руководитель: Попов Игорь Юрьевич, д.ф.-м.н.
}

%\institute[ИТМО]{НОЦМ, Университет ИТМО}

%\title{From Planck data to Planck era: \\ Observational tests of Holographic Cosmology} 

%\author{Elizabeth Gould} \affiliation{Perimeter Institute for Theoretical Physics, 31 Caroline St. N., Waterloo, ON, N2L 2Y5, Canada} \affiliation{Department of Physics and Astronomy, University of Waterloo, Waterloo, ON, N2L 3G1, Canada}
 
\date{\today}

\maketitle
%\makeindex




\chapter{Introduction}
%input{Introduction}
\chapter{Error}
%input{Error}
\chapter{The IVP}
%input{Notebook_to_build_graphics}
\chapter{Fast Oscillation Modeling}
%input{Notebook_Fast_t_IVP}
%input{Notebook_IVP_fast_t_MC_tests}
\chapter{Slow Oscillation Modeling}
%input{Notebook_Slow_k_IVP}
%input{Notebook_to_build_graphics}
\chapter{bvp matching}
%input{Notebook_BVP_Plots}
\chapter{complete solution}
%input{Notebook_BVP_to_end}
%input{Notebook_mu_vs_I}
\chapter{remodeling k and m}
%input{Notebook_BVP_k_M_structure}
\chapter{Conclusion}
%input{Conclusion}

%input{}

\lstset{style=codejupytr}

\begin{lstlisting}
test
\end{lstlisting}


% The \appendix statement indicates the beginning of the appendices.
\appendix

% Add a title page before the appendices and a line in the Table of Contents
%\chapter*{APPENDICES}
%\addcontentsline{toc}{chapter}{APPENDICES}


%\lstinputlisting{FinalizedFiles/eelib/Translation/loop.py}
\lstset{style=codepy}

\chapter{init.py}
\begin{lstlisting}[language=Python]
# __init__.py

# Author: Elizabeth Gould
# Builds the library. Note that all functions, constants, and classes are added twice.
# First they are imported, then assigned.

# Helper functions and functions not fully implemented don't appear here, as they are not intended to be 
# used from the library. If they are needed, they must be called by eelib.filename.function_name instead 
# of eelib.function_name. Note that the format eelib.filename.function_name always works.

from eelib.consts import h_pl, hbar, m_e, e_e, c_l, pi, rtol, atol, MU_min, MU_max, DK_min, DK_max
from eelib.consts import kFAu, R_max, R_min, B_min, B_max, I_e, phi0inv, pppterm
from eelib.deriv_functions import psi_deriv, psi_deriv_old, psi_deriv_0, psi_deriv_full
from eelib.table_scripts import clean_table, find_ind_var, clean_table_MC

__all__ = ['h_pl', 'hbar', 'm_e', 'e_e', 'c_l', 'pi', 'kFAu', 'R_max', 'R_min', 'B_min', 'B_max', 'I_e', 
           'phi0inv', 'pppterm', 'rtol', 'atol', 'MU_min', 'MU_max', 'DK_min', 'DK_max', 
           'psi_deriv', 'psi_deriv_old', 'psi_deriv_0', 'psi_deriv_full', 
           'event1', 'event2', 'deriv_amp', 'deriv_phase', 'deriv_real']
\end{lstlisting}

%% document class
\documentclass[A4,12pt,oneside]{report}
%A4,12pt,titlepage,oneside,final

% list of packages
\usepackage[left=2cm,right=2cm,top=2cm,bottom=2cm]{geometry} % document size

\usepackage[utf8]{inputenc}
\usepackage[T2A]{fontenc} %russian language encoding

\usepackage{textcomp}
\usepackage{lmodern}
\usepackage{amsmath}
\usepackage{hyperref}
\usepackage{xcolor}
\usepackage{graphicx}
\usepackage{listings}
\usepackage{listingsutf8}
\usepackage{dcolumn}   % needed for some tables
\usepackage{bm}        
\usepackage{amssymb} 
\usepackage{fancyvrb}
\usepackage{fvextra}
%\usepackage{minted}

%\setcounter{secnumdepth}{3}
%\setcounter{tocdepth}{3}
%\usepackage{eso-pic}
\usepackage{array}
%\usepackage{rotating}
\usepackage{multicol}
\usepackage{setspace}

%russian language support
\usepackage[english,russian]{babel}
%starts at the last language on the list
%select language switches the language


\lstset{breaklines=true}

% The following is a list of styles for use with listings

\lstdefinestyle{codepy}{
    backgroundcolor=\color{white},   
    commentstyle=\color{green!50!black},
    keywordstyle=\color{blue},
    numberstyle=\tiny\color{gray},
    stringstyle=\color{red!50!black},
    basicstyle=\ttfamily\footnotesize,
    breakatwhitespace=false,         
    breaklines=true,                 
    captionpos=b,                    
    keepspaces=true,                 
    numbers=left,                    
    numbersep=5pt,                  
    showspaces=false,                
    showstringspaces=false,
    showtabs=false,                  
    tabsize=2,
	texcl=true,     %<---- added
	escapeinside=@@,
	%escapebegin=\begin{russian}\commentfont,
	%escapeend=\end{russian},
}

\makeatletter


\graphicspath{{figs/}} % Directory in which figures are stored


%
% list of custom commands
\newcommand{\be}{\begin{equation}}
\newcommand{\ee}{\end{equation}}
\newcommand{\bea}{\begin{eqnarray}}
\newcommand{\eea}{\end{eqnarray}}
\renewcommand{\(}{\left(}
\renewcommand{\)}{\right)}

\newcommand{\rut}{\selectlanguage{russian}}
\newcommand{\ent}{\selectlanguage{english}}


%\newcommand{\be}{\begin{equation}}
%\rnc{\d}{\mathrm{d}}
%\nc{\D}{\partial}
%\nc{\nn}{\nonumber}
%\rnc{\inf}{\infty}
%\nc{\vphi}{\varphi}
%\nc{\half}{\frac{1}{2}}


% color commands are used to highlight text temporarily, for editing
\newcommand{\nadd}[1]{\textcolor{magenta}{#1}}
\newcommand{\ncomm}[1]{\textcolor{magenta}{\bf{[#1]}}}
\renewcommand{\Re}{\mathrm{Re}}
\renewcommand{\Im}{\mathrm{Im}}




\makeatother

\begin{document}

\lstset{style=codepy}
\begin{lstlisting}[language = Python]
# \_\_init\_\_.py

# Автор: Элизабет Гоулд
# Дата последнего изменения: 19.03.2026

# Создает библиотеку. Все функции, константы и классы добавляются дважды.
# Сначала они импортируются, а затем назначаются.

# Вспомогательные функции и функции, реализованные не полностью, здесь не 
# отображаются, поскольку они не предназначены для использования из библиотеки.
# Если они необходимы, они должны вызываться с помощью
# eelib.filename.function\_name вместо eelib.function\_name.
# Формат eelib.filename.function\_name всегда работает.

from eelib.consts import h_pl, hbar, m_e, e_e, c_l, pi, rtol, atol, MU_min, MU_max, DK_min, DK_max
from eelib.consts import kFAu, R_max, R_min, B_min, B_max, I_e, phi0inv, pppterm
from eelib.deriv_functions import psi_deriv, psi_deriv_old, psi_deriv_0, psi_deriv_full
from eelib.events import event1, event2, deriv_amp, deriv_phase, deriv_real

from eelib.fitted_functions import sin_fit
from eelib.fitted_functions import fit_sin

from eelib.k_M_models_ivp import pred_fast_t, pred_fast_k, pred_slow_k, pred_slow_k_v3, pred_fast_k_true, pred_slow_k_v2
from eelib.bvp_rootfinder_functions import function_wrapper_bvp, k_calc_0, M_calc_0
from eelib.bvp_test_script import test_match

from eelib.loop import loop
from eelib.deriv_grid import deriv_grid
from eelib.grid_fast_osc import grid_fast_osc
from eelib.grid_slow_osc import grid_slow_osc
from eelib.grid_bvp import grid_BVP

from eelib.table_scripts import clean_table, find_ind_var, clean_table_MC

__all__ = ['h_pl', 'hbar', 'm_e', 'e_e', 'c_l', 'pi', 'kFAu', 'R_max', 'R_min', 'B_min', 'B_max', 'I_e', 
           'phi0inv', 'pppterm', 'rtol', 'atol', 'MU_min', 'MU_max', 'DK_min', 'DK_max', 
           'psi_deriv', 'psi_deriv_old', 'psi_deriv_0', 'psi_deriv_full', 
           'event1', 'event2', 'deriv_amp', 'deriv_phase', 'deriv_real', 
           'sin_fit', 'fit_sin', 'clean_table', 'find_ind_var', 'clean_table_MC',
           'pred_fast_t', 'pred_fast_k', 'pred_slow_k', 'pred_slow_k_v3', 'pred_fast_k_true',
           'pred_slow_k_v2', 'function_wrapper_bvp', 'test_match', 'k_calc_0', 'M_calc_0',
           'loop', 'deriv_grid', 'grid_fast_osc', 'grid_slow_osc', 'grid_BVP']
\end{lstlisting}

\end{document}

%\input{consts}


\end{document}
